% Generated draft: requires mathematical and source review.
\documentclass[11pt]{article}
\usepackage[margin=25mm]{geometry}
\usepackage{fontspec}
\setmainfont{Latin Modern Roman}
\usepackage{amsmath,amssymb,mathtools,unicode-math}
\setmathfont{Latin Modern Math}
\usepackage{enumitem,needspace,xurl,hyperref}
\setlength{\parindent}{0pt}
\setlength{\parskip}{5pt}
\setlength{\emergencystretch}{3em}
\newcommand{\N}{\mathbb{N}}
\newcommand{\Z}{\mathbb{Z}}
\newcommand{\Q}{\mathbb{Q}}
\newcommand{\R}{\mathbb{R}}
\newcommand{\C}{\mathbb{C}}
\newcommand{\F}{\mathbb{F}}
\newcommand{\A}{\mathbb{A}}
\newcommand{\PP}{\mathbb{P}}
\newcommand{\E}{\mathbb{E}}
\newcommand{\eps}{\varepsilon}
\newcommand{\dd}{\,\mathrm{d}}
\begin{document}
\section*{UnsolvedMath ID: 11000126}
\textbf{Source problem code:} catalogue entry
\subsection{Definitions and mathematical statement}
Let \(\Gamma_\infty\) be the stable mapping class group under the source's genus-stabilization convention, and \(B\Gamma_\infty\) its classifying space. In the cohomological reading of the extracted notation, the problem asks for a direct description of some particular simple torsion classes in
\[
 H^*(B\Gamma_\infty;\Z).
\]
A torsion class here is a nonzero \(x\in H^q(B\Gamma_\infty;\Z)\) with \(mx=0\) for some integer \(m>0\). The target is explicit examples with a direct construction and justified torsion/nonvanishing, not merely an abstract existence result or a computation of rational classes. The source does not specify degrees, orders, a simplicity criterion or the allowed meaning of direct.

It additionally asks to study the relation of higher Reidemeister torsion from [10] to the classes \(\kappa_i\), noting that both [10] and [14] use parametrized Morse theory, and asks whether ``modular higher Reidemeister torsion'' could relate to this problem. The usual \(\kappa_i\) are Miller--Morita--Mumford surface-bundle classes, expressed by fiber integration of the \((i+1)\)-st power of the vertical Euler class, in degree \(2i\); the intended relative/boundary version is not supplied. No definition of modular torsion or comparison identity is printed.

The following context concerns a different stable group: for \(F_g\) free on \(g\) generators, extend each automorphism by fixing the new generator to form \(A_\infty=\varinjlim_g\operatorname{Aut}(F_g)\). This is the stable group denoted \(\operatorname{Aut}(F_\infty)\) in the source, not automatically the full automorphism group of a countably generated free group. The source reports a perfect index-two subgroup, forms \(BA_\infty^+\), and attributes to A. Hatcher [8] a splitting in which the stable sphere infinite-loop space, printed \(\Omega^\infty S^\infty\), is a homotopy direct factor of \(\Z\times BA_\infty^+\). This background is not strengthened to an equivalence or used as a solution to the mapping-class torsion problem.
\subsection{Short English statement}
Construct explicit simple integral torsion classes in stable mapping-class cohomology, and investigate whether a modular form of higher Reidemeister torsion explains them or their relation to the \(\kappa_i\) classes.
\subsection{Sources}
\begin{itemize}
\item \emph{Problems on Mapping Class Groups and Related Topics}, edited by Farb, Problem 3.3 and following discussion, supplied PDF page 104. \url{https://www.math.uchicago.edu/~farb/papers/mcgbook.pdf}. Recorded access 2026-07-29; extraction pdf-font-structure.
\item Source [10] for higher Reidemeister torsion, [14] for the other parametrized-Morse-theory construction, and A. Hatcher [8] for the stable free-group splitting. Full bibliography entries are absent. The supplied summary attributes torsion existence to Galatius and higher-Reidemeister-torsion expressions to Igusa without complete citations.
\item Local UnsolvedMath registry, \texttt{data/unsolvedmath/problems.json}, permanent ID \(11000126\), exact source code \texttt{catalogue entry}; supplied OPEN-TRIAGE summary reports no located direct elementary construction.
\end{itemize}
\end{document}
